\documentclass[12pt,letterpaper]{article}
\usepackage[utf8]{inputenc}
\usepackage[spanish]{babel}
\usepackage[margin=2.5cm]{geometry}
\usepackage{amsmath}
\usepackage{amsfonts}
\usepackage{amssymb}

\begin{document}

\begin{center}
    \textbf{\Large Examen 1: Unidad I (Clave de Respuestas)}\\
    \vspace{0.2cm}
    \textbf{Física para Ciencias de la Salud (005-1713)}\\
\end{center}

\vspace{0.5cm}

\begin{enumerate}

    \item \textbf{Ángulo plano:} Un paciente en rehabilitación realiza un movimiento de flexión con su brazo describiendo un ángulo de $75^\circ$. Exprese la medida de este ángulo en radianes. \\
    \textbf{Solución:} \\
    Para convertir de grados a radianes, multiplicamos por el factor de conversión $\frac{\pi \text{ rad}}{180^\circ}$:
    \[ \theta = 75^\circ \times \left( \frac{\pi}{180^\circ} \right) = \frac{75\pi}{180} \text{ rad} = \frac{5\pi}{12} \text{ rad} \approx 1.309 \text{ rad} \]

    \item \textbf{Sistemas de coordenadas:} En una placa de rayos X digitalizada 2D, el punto de origen de un hueso (articulación A) está en las coordenadas $(1, 2)$ cm y el extremo B en $(7, 10)$ cm. Calcule la longitud real del hueso. \\
    \textbf{Solución:} \\
    Aplicamos la fórmula de distancia entre dos puntos:
    \[ d = \sqrt{(x_2 - x_1)^2 + (y_2 - y_1)^2} \]
    \[ d = \sqrt{(7 - 1)^2 + (10 - 2)^2} = \sqrt{6^2 + 8^2} = \sqrt{36 + 64} = \sqrt{100} = 10 \text{ cm} \]
    
    \item \textbf{Vectores:} Sobre la vértebra L5 de una persona actúan dos fuerzas debido a la tensión muscular: $\vec{F}_1 = (45\hat{i} + 25\hat{j})$ N y $\vec{F}_2 = (-15\hat{i} + 35\hat{j})$ N. Encuentre el vector de la fuerza resultante que soporta la vértebra. \\
    \textbf{Solución:} \\
    Sumamos componente a componente:
    \[ \vec{R} = \vec{F}_1 + \vec{F}_2 = (45 - 15)\hat{i} + (25 + 35)\hat{j} = 30\hat{i} + 60\hat{j} \text{ N} \]
    
    \item \textbf{Potencias de diez:} Un glóbulo blanco (leucocito) humano tiene un diámetro aproximado de $0.0000125$ metros. Exprese esta medida utilizando notación científica (potencias de diez) y el prefijo adecuado del Sistema Internacional. \\
    \textbf{Solución:} \\
    Movemos la coma decimal 5 lugares a la derecha para la notación científica estándar ($1.25 \times 10^{-5}$ m). Usando un prefijo común (micro, $10^{-6}$):
    \[ d = 12.5 \times 10^{-6} \text{ m} = 12.5 \, \mu\text{m} \]

    \item \textbf{Sistemas de unidades:} En un análisis de laboratorio, se determina que la densidad del plasma sanguíneo es de $1.03 \, \text{g/cm}^3$. Convierta este valor a las unidades estándar del Sistema Internacional ($\text{kg/m}^3$). \\
    \textbf{Solución:} \\
    \[ \rho = 1.03 \, \frac{\text{g}}{\text{cm}^3} \times \left( \frac{1 \text{ kg}}{1000 \text{ g}} \right) \times \left( \frac{10^6 \text{ cm}^3}{1 \text{ m}^3} \right) = 1030 \, \frac{\text{kg}}{\text{m}^3} \]

\end{enumerate}

\end{document}